\begin{frame}{동역학: 위치에서 운동으로 --- 진자의 방정식}
  기하학은 ``로봇이 \emph{어딘가}에 있다''만 말할 뿐, ``힘을 주면
  \emph{어떻게} 움직이는가''는 못 함. 후자가 \kb{동역학}(dynamics).
  MDP 전이함수 $P(s' \mid s, a)$ 가 ``시뮬레이션 한 스텝''이라면,
  그 한 스텝의 \emph{내용}이 바로 동역학 방정식.
  \vskip0.5em
  가장 단순한 예 \kb{진자}: $\theta$ 는 세로(위)에서 잰 각도
  ($0$ 이 위쪽 수직 = \textbf{불안정한} 균형점, $\pm\pi$ 가 아래 매달린
  \textbf{안정}점). 균일한 막대($m, l$, $I = \tfrac{1}{3}ml^2$)에
  뉴턴-오일러 방정식(중력 토크 $\tfrac{1}{2}mgl\sin\theta$):
  \[
    \tfrac{1}{3}ml^2 \,\ddot\theta = \tfrac{1}{2} mgl \sin\theta + u
    \;\;\Longrightarrow\;\;
    \ddot\theta = \frac{3g}{2l}\sin\theta + \frac{3}{ml^2}\,u
  \]
  $u$ = 우리가 주는 토크(\textbf{행동!}).
  \begin{itemize}
    \item (1) $\sin\theta$ 항 = 위쪽 수직에서 \textbf{떨어지려 하는 힘} ---
      0(위) 근처에서 $\ddot\theta > 0$, 더 벌어진 쪽으로 가속.
      즉 0은 \emph{불안정한} 균형(진자를 세우기 어려운 이유).
    \item (2) $u$ 의 계수 $3/(ml^2)$: 무거울수록($m$ ↑), 길수록($l$ ↑)
      같은 토크가 작은 가속을 준다 --- \textbf{관성의 법칙} 그 자체.
  \end{itemize}
\end{frame}
